\section{September 11, 2026}

We continue the discussion of stability from the previous lecture. The
main goals are to linearize an autonomous ODE near an equilibrium, to
study the resulting system using the matrix exponential, and to extend
the solution formula to time-dependent linear systems by means of the
Dyson series.

\subsection{Linearization near an equilibrium}

Consider the autonomous nonlinear equation
\begin{equation}\label{eq:nonlinear-autonomous-sep-11}
  x'=f(x)
\end{equation}
and let \(x_*\) be an equilibrium, so that \(f(x_*)=0\). Introduce the
perturbation
\[
  y=x-x_*.
\]
Since \(x_*\) is constant, \(y'=x'\). If \(f\) is continuously
differentiable, its expansion near \(x_*\) is
\[
  f(x_*+y)=f(x_*)+f_x(x_*)y+r(y)
           =Ay+r(y),
  \qquad
  A=f_x(x_*),
\]
where
\[
  \frac{\lVert r(y)\rVert}{\lVert y\rVert}\longrightarrow0
  \qquad\text{as }y\to0.
\]
Thus the exact perturbation equation is
\begin{equation}\label{eq:perturbation-equation-sep-11}
  y'=Ay+r(y),
\end{equation}
and its linearization is
\begin{equation}\label{eq:linearized-equation-sep-11}
  y'=Ay.
\end{equation}
The equilibrium \(x_*\) corresponds to the zero solution of the
perturbation equation. Hence studying whether \(x(t)\) remains close to
\(x_*\) begins with studying whether solutions of
\eqref{eq:linearized-equation-sep-11} remain close to zero.

\subsection{Autonomous linear systems and the matrix exponential}

Consider the initial-value problem
\begin{equation}\label{eq:linear-autonomous-sep-11}
  y'=Ay,
  \qquad
  y(t_0)=y_0,
\end{equation}
where \(A\in\bb{R}^{n\times n}\) is constant. Repeated differentiation
gives
\[
  y^{(k)}(t_0)=A^ky_0.
\]
The Taylor series about \(t_0\) therefore gives
\begin{align*}
  y(t)
  &=y_0+(t-t_0)Ay_0
    +\frac{(t-t_0)^2}{2!}A^2y_0+\cdots\\
  &=\left(\sum_{k=0}^{\infty}
    \frac{(t-t_0)^kA^k}{k!}\right)y_0.
\end{align*}

\begin{definition}[Matrix exponential]
  For a square matrix \(B\), its exponential is
  \begin{equation}\label{eq:matrix-exponential-definition-sep-11}
    e^B=\sum_{k=0}^{\infty}\frac{B^k}{k!}.
  \end{equation}
\end{definition}

The series converges in every matrix norm. With this notation, the
solution of \eqref{eq:linear-autonomous-sep-11} is
\begin{equation}\label{eq:linear-solution-sep-11}
  y(t)=e^{(t-t_0)A}y_0.
\end{equation}
Indeed,
\[
  \frac{d}{dt}e^{(t-t_0)A}
  =Ae^{(t-t_0)A}=e^{(t-t_0)A}A,
  \qquad
  e^{0A}=I.
\]

As observed in the previous lecture, the zero solution of
\eqref{eq:linear-autonomous-sep-11} is stable exactly when
\begin{equation}\label{eq:bounded-semigroup-sep-11}
  \sup_{t\geq0}\lVert e^{tA}\rVert<\infty,
\end{equation}
and it is asymptotically stable exactly when
\begin{equation}\label{eq:decaying-semigroup-sep-11}
  \lVert e^{tA}\rVert\longrightarrow0
  \qquad\text{as }t\to\infty.
\end{equation}
These statements follow directly from
\[
  y(t;t_0,y_0)-y(t;t_0,z_0)
  =e^{(t-t_0)A}(y_0-z_0).
\]
Linearity also implies that stability of the zero solution is
equivalent to stability of every solution.

\subsection{Time-dependent linear systems and the Dyson series}

Now consider the nonautonomous linear equation
\begin{equation}\label{eq:nonautonomous-linear-sep-11}
  y'=A(t)y,
  \qquad
  y(t_0)=y_0,
\end{equation}
where \(A(t)\) is continuous. Its integral equation is
\begin{equation}\label{eq:linear-integral-equation-sep-11}
  y(t)=y_0+\int_{t_0}^t A(t_1)y(t_1)\,dt_1.
\end{equation}
At an intermediate time \(t_1\), the same equation reads
\[
  y(t_1)=y_0+\int_{t_0}^{t_1}A(t_2)y(t_2)\,dt_2.
\]
Substituting this formula into
\eqref{eq:linear-integral-equation-sep-11} gives
\begin{align*}
  y(t)
  ={}&y_0+\left(\int_{t_0}^t A(t_1)\,dt_1\right)y_0\\
  &+\int_{t_0}^t\int_{t_0}^{t_1}
    A(t_1)A(t_2)y(t_2)\,dt_2\,dt_1.
\end{align*}
Continuing this substitution produces the \emph{Dyson series}
\begin{equation}\label{eq:dyson-series-sep-11}
  y(t)=U(t,t_0)y_0,
\end{equation}
where the evolution operator is
\begin{align}
  U(t,t_0)
  ={}&I+\sum_{k=1}^{\infty}
  \int_{t_0}^t\int_{t_0}^{t_1}\cdots
  \int_{t_0}^{t_{k-1}}
  A(t_1)A(t_2)\cdots A(t_k)
  \,dt_k\cdots dt_2\,dt_1.
  \label{eq:evolution-operator-dyson-sep-11}
\end{align}
If \(\lVert A(s)\rVert\leq M\) on \([t_0,t]\), the norm of the
\(k\)-th term is at most
\[
  \frac{M^k(t-t_0)^k}{k!},
\]
which proves convergence of the series. Termwise differentiation shows
that
\[
  \frac{\partial}{\partial t}U(t,t_0)=A(t)U(t,t_0),
  \qquad
  U(t_0,t_0)=I.
\]

Because matrices evaluated at different times need not commute, the
order \(t_1\geq t_2\geq\cdots\geq t_k\) in
\eqref{eq:evolution-operator-dyson-sep-11} is essential. A common
shorthand is the \emph{time-ordered exponential}
\begin{equation}\label{eq:time-ordered-exponential-sep-11}
  U(t,t_0)
  =\mathcal{T}\exp\left(\int_{t_0}^t A(s)\,ds\right).
\end{equation}
If \(A(s)A(r)=A(r)A(s)\) for every \(r,s\), time ordering is
unnecessary and
\[
  U(t,t_0)=\exp\left(\int_{t_0}^t A(s)\,ds\right).
\]
In particular, when \(A(t)\equiv A\), the ordered integration region
has volume \((t-t_0)^k/k!\), and the Dyson series reduces to
\(e^{(t-t_0)A}\).

\subsection{Properties of the matrix exponential}

The power-series definition immediately gives several useful
identities.

\begin{proposition}\label{prop:matrix-exponential-properties-sep-11}
  Let \(A\) and \(B\) be square matrices of the same size.
  \begin{enumerate}[label=\textup{(\roman*)}]
    \item If \(M\) is invertible, then
      \[
        e^{MAM^{-1}}=Me^AM^{-1}.
      \]
    \item If \(AB=BA\), then
      \[
        e^{A+B}=e^Ae^B=e^Be^A.
      \]
    \item For a block-diagonal matrix,
      \[
        e^{\operatorname{diag}(A_1,\ldots,A_m)}
        =\operatorname{diag}(e^{A_1},\ldots,e^{A_m}).
      \]
  \end{enumerate}
\end{proposition}

\begin{proof}
  For the first identity, observe that
  \((MAM^{-1})^k=MA^kM^{-1}\) and apply
  \eqref{eq:matrix-exponential-definition-sep-11}. The third identity
  follows similarly because powers preserve block-diagonal structure.
  When \(A\) and \(B\) commute, the binomial theorem is valid for
  \((A+B)^k\). Substitution into the exponential series and collection
  of like terms gives the second identity.
\end{proof}

\subsection{Jordan form and linear stability}

The preceding identities reduce the stability analysis of \(y'=Ay\)
to the analysis of individual Jordan blocks. Over \(\bb{C}\), write
\begin{equation}\label{eq:jordan-decomposition-sep-11}
  A=PJP^{-1},
  \qquad
  J=\operatorname{diag}(J_1,\ldots,J_m).
\end{equation}
If \(J_k\) is a Jordan block of size \(q_k\) associated with the
eigenvalue \(\lambda_k\), then
\[
  J_k=\lambda_k I+N_k,
  \qquad
  N_k^{q_k}=0.
\]
Since \(\lambda_kI\) commutes with \(N_k\),
\begin{equation}\label{eq:jordan-block-exponential-sep-11}
  e^{tJ_k}
  =e^{t\lambda_k}
   \left(
     I+tN_k+\frac{t^2N_k^2}{2!}+\cdots
     +\frac{t^{q_k-1}N_k^{q_k-1}}{(q_k-1)!}
   \right).
\end{equation}
Thus each Jordan block contributes an exponential factor multiplied by
a polynomial in \(t\).

\begin{theorem}[Stability criterion for a linear system]
  For \(y'=Ay\), the following statements hold.
  \begin{enumerate}[label=\textup{(\roman*)}]
    \item The zero solution is asymptotically stable if and only if
      \[
        \operatorname{Re}\lambda<0
        \qquad\text{for every }\lambda\in\sigma(A).
      \]
    \item The zero solution is stable if and only if every eigenvalue
      satisfies \(\operatorname{Re}\lambda\leq0\) and every eigenvalue
      on the imaginary axis is semisimple, meaning that all of its
      Jordan blocks have size one.
  \end{enumerate}
\end{theorem}

\begin{proof}
  By Proposition~\ref{prop:matrix-exponential-properties-sep-11} and
  \eqref{eq:jordan-decomposition-sep-11},
  \[
    e^{tA}=Pe^{tJ}P^{-1}
    =P\operatorname{diag}(e^{tJ_1},\ldots,e^{tJ_m})P^{-1}.
  \]
  If \(\operatorname{Re}\lambda_k<0\), the exponential decay in
  \eqref{eq:jordan-block-exponential-sep-11} dominates every
  polynomial factor, so the block tends to zero. If
  \(\operatorname{Re}\lambda_k>0\), the block has an exponentially
  growing solution. When \(\operatorname{Re}\lambda_k=0\), the scalar
  exponential has constant magnitude. The block is therefore bounded
  precisely when its nilpotent part is zero, or equivalently when its
  size is one. Combining the blocks with
  \eqref{eq:bounded-semigroup-sep-11} and
  \eqref{eq:decaying-semigroup-sep-11} proves both assertions.
\end{proof}

\begin{example}[A defective eigenvalue on the imaginary axis]
  Let
  \[
    A=\begin{bmatrix}0&1\\0&0\end{bmatrix}.
  \]
  The only eigenvalue is zero, but its Jordan block has size two. Since
  \(A^2=0\),
  \[
    e^{tA}=I+tA
    =\begin{bmatrix}1&t\\0&1\end{bmatrix}.
  \]
  Hence
  \[
    y_2(t)=y_2(0),
    \qquad
    y_1(t)=y_1(0)+t\,y_2(0).
  \]
  Arbitrarily small initial data with \(y_2(0)\neq0\) lead to linear
  growth, so the origin is not stable. This illustrates why
  nonpositive real parts alone are insufficient for stability.
\end{example}

\enlargethispage{2\baselineskip}
For the nonlinear perturbation equation
\eqref{eq:perturbation-equation-sep-11}, the linearization principle
gives two decisive cases. If the spectrum of \(f_x(x_*)\) lies in the
open left half-plane, then \(x_*\) is locally asymptotically stable. If
the spectrum contains an eigenvalue with positive real part, then
\(x_*\) is unstable. If it meets the imaginary axis but not the open
right half-plane, linearization alone may be inconclusive; the
nonlinear remainder \(r(y)\) can determine the behavior.
